\sectionkicker{Source-backed technical notes}
\subsection*{Agent Safety \& Security — 能力、権限、接続面を分けて考える: Technical Notes}
本欄は記事本文で圧縮した一次資料上の情報を、比較・再検証しやすい形で整理したものである。時系列、確認済みの主張、留意点、一次資料URLを掲載する。
\medskip
\subsection*{Theme at a glance}
\addcontentsline{toc}{subsection}{Theme at a glance}
\begin{center}
\footnotesize
\begin{tabularx}{\linewidth}{@{}p{0.20\linewidth}p{0.10\linewidth}p{0.14\linewidth}X@{}}
\toprule
資料 & 位置づけ & 種別 & 時系列 \\
\midrule
Capability Gates Are Not Authorization: Confused-Deputy Failures in LLM Agent Frameworks & 主要資料 & 論文 & 2026-06-27 (論文\_RELEASE) \\
Adaptive Evaluation of Out-of-Band Defenses Against Prompt Injection in LLM Agents & 主要資料 & 論文 & 2026-06-25 (論文\_RELEASE) \\
Understanding and Evaluating Claw-like Agent Security Through a Computer-Systems Lens & 主要資料 & 論文 & 2026-06-29 (論文\_RELEASE) \\
ShareLock: A Stealthy Multi-Tool Threshold Poisoning Attack Against MCP & 補足資料 & 論文 & 2026-06-25 (論文\_RELEASE) \\
Deployment Simulation & 補足資料 & 安全性事象 & 2026-06-16 (SAFETY\_METHOD\_RELEASE) \\
\bottomrule
\end{tabularx}
\end{center}
\normalsize

\begin{technicalnote}{Capability Gates Are Not Authorization: Confused-Deputy Failures in LLM Agent Frameworks}{主要資料}
\begin{tabularx}{\linewidth}{@{}>{\bfseries}p{0.22\linewidth}X@{}}
組織 & Paper authors \\
種別 & 論文 \\
時系列 & 2026-06-27 (論文\_RELEASE) \\
\end{tabularx}
\smallskip
{\bfseries 一次資料から整理したtechnical points}
\begin{itemize}[leftmargin=1.5em,itemsep=0.35em]
\item \textbf{Author claim}: The authors distinguish tool capability exposure from per-call authorization and demonstrate confused-deputy risks in common agent frameworks.
\end{itemize}
{\bfseries 読む際の境界}
\begin{itemize}[leftmargin=1.5em,itemsep=0.35em]
\item \textbf{分析上の留意点}: The evidence is the authors' arXiv paper/model card and reported evaluations; results are not independently reproduced in this Evidence review.
\end{itemize}
{\bfseries 一次資料}
\begingroup\sloppy
\begin{itemize}[leftmargin=1.5em,itemsep=0.25em]
\item {\scriptsize\url{http://arxiv.org/abs/2606.28679v1}}
\end{itemize}
\endgroup
\end{technicalnote}
\medskip

\begin{technicalnote}{Adaptive Evaluation of Out-of-Band Defenses Against Prompt Injection in LLM Agents}{主要資料}
\begin{tabularx}{\linewidth}{@{}>{\bfseries}p{0.22\linewidth}X@{}}
組織 & Paper authors \\
種別 & 論文 \\
時系列 & 2026-06-25 (論文\_RELEASE) \\
\end{tabularx}
\smallskip
{\bfseries 一次資料から整理したtechnical points}
\begin{itemize}[leftmargin=1.5em,itemsep=0.35em]
\item \textbf{Author claim}: The authors argue prompt-injection defenses must be evaluated with adaptive, defense-aware attacks rather than a fixed attack suite.
\end{itemize}
{\bfseries 読む際の境界}
\begin{itemize}[leftmargin=1.5em,itemsep=0.35em]
\item \textbf{分析上の留意点}: The evidence is the authors' arXiv paper/model card and reported evaluations; results are not independently reproduced in this Evidence review.
\end{itemize}
{\bfseries 一次資料}
\begingroup\sloppy
\begin{itemize}[leftmargin=1.5em,itemsep=0.25em]
\item {\scriptsize\url{http://arxiv.org/abs/2606.26479v1}}
\end{itemize}
\endgroup
\end{technicalnote}
\medskip

\begin{technicalnote}{Understanding and Evaluating Claw-like Agent Security Through a Computer-Systems Lens}{主要資料}
\begin{tabularx}{\linewidth}{@{}>{\bfseries}p{0.22\linewidth}X@{}}
組織 & Paper authors \\
種別 & 論文 \\
時系列 & 2026-06-29 (論文\_RELEASE) \\
\end{tabularx}
\smallskip
{\bfseries 一次資料から整理したtechnical points}
\begin{itemize}[leftmargin=1.5em,itemsep=0.35em]
\item \textbf{Author claim}: The authors frame persistent Claw-like agents as computer systems and introduce SafeClawArena for cross-component security evaluation.
\end{itemize}
{\bfseries 読む際の境界}
\begin{itemize}[leftmargin=1.5em,itemsep=0.35em]
\item \textbf{分析上の留意点}: The evidence is the authors' arXiv paper/model card and reported evaluations; results are not independently reproduced in this Evidence review.
\end{itemize}
{\bfseries 一次資料}
\begingroup\sloppy
\begin{itemize}[leftmargin=1.5em,itemsep=0.25em]
\item {\scriptsize\url{http://arxiv.org/abs/2606.30755v1}}
\end{itemize}
\endgroup
\end{technicalnote}
\medskip

\begin{technicalnote}{ShareLock: A Stealthy Multi-Tool Threshold Poisoning Attack Against MCP}{補足資料}
\begin{tabularx}{\linewidth}{@{}>{\bfseries}p{0.22\linewidth}X@{}}
組織 & Paper authors \\
種別 & 論文 \\
時系列 & 2026-06-25 (論文\_RELEASE) \\
\end{tabularx}
\smallskip
{\bfseries 一次資料から整理したtechnical points}
\begin{itemize}[leftmargin=1.5em,itemsep=0.35em]
\item \textbf{Author claim}: The authors introduce ShareLock, a threshold-based multi-tool poisoning attack for MCP ecosystems.
\end{itemize}
{\bfseries 読む際の境界}
\begin{itemize}[leftmargin=1.5em,itemsep=0.35em]
\item \textbf{分析上の留意点}: The evidence is the authors' arXiv paper/model card and reported evaluations; results are not independently reproduced in this Evidence review.
\end{itemize}
{\bfseries 一次資料}
\begingroup\sloppy
\begin{itemize}[leftmargin=1.5em,itemsep=0.25em]
\item {\scriptsize\url{http://arxiv.org/abs/2606.27027v1}}
\end{itemize}
\endgroup
\end{technicalnote}
\medskip

\begin{technicalnote}{Deployment Simulation}{補足資料}
\begin{tabularx}{\linewidth}{@{}>{\bfseries}p{0.22\linewidth}X@{}}
組織 & OpenAI \\
種別 & 安全性事象 \\
時系列 & 2026-06-16 (SAFETY\_METHOD\_RELEASE) \\
\end{tabularx}
\smallskip
{\bfseries 一次資料から整理したtechnical points}
\begin{itemize}[leftmargin=1.5em,itemsep=0.35em]
\item \textbf{Vendor claim}: OpenAI introduced Deployment Simulation, which replays prior conversations with a candidate model and simulates tool behavior to estimate pre-deployment behavior in more realistic contexts.
\end{itemize}
{\bfseries 読む際の境界}
\begin{itemize}[leftmargin=1.5em,itemsep=0.35em]
\item \textbf{分析上の留意点}: The source is first-party OpenAI material. Capability, benchmark, efficiency, or performance statements remain vendor claims unless explicitly described as externally reproduced.
\end{itemize}
{\bfseries 一次資料}
\begingroup\sloppy
\begin{itemize}[leftmargin=1.5em,itemsep=0.25em]
\item {\scriptsize\url{https://openai.com/index/deployment-simulation}}
\end{itemize}
\endgroup
\end{technicalnote}
\medskip
